
\section{Method}

\subsection{Overview}

Our implementation wraps FlashAttention-3 with three layers of safety:

\begin{enumerate}
\item \textbf{Import-time detection.} At module load, we attempt to import
  \texttt{flash\_attn\_interface} or \texttt{flash\_attn}. If neither is available, we set
  \texttt{FLASH\_ATTN\_AVAILABLE = False}.

\item \textbf{Runtime FP8 detection.} At each forward, we check that:
\end{enumerate}
\begin{itemize}
\item \texttt{FLASH\_ATTN\_AVAILABLE} is True
\item \texttt{self.use\_fp8} is True
\item \texttt{q.is\_cuda} is True
\item \texttt{torch.float8\_e4m3fn} is available
  
  If all hold, we cast Q/K/V to FP8 and call \texttt{flash\_attn\_func}.

\end{itemize}
\begin{enumerate}
\item \textbf{SDPA fallback.} If any check fails, we call
  \texttt{F.scaled\_dot\_product\_attention} instead.

\end{enumerate}
\begin{lstlisting}
              +-----------------------------+
              |  Input Q, K, V              |
              +--------------+--------------+
                             |
                              
              +-----------------------------+
              |  FLASH\_ATTN\_AVAILABLE ?     |
              +------+--------------+-------+
                     | No           | Yes
                     |              |
                     |               
                     |    +-----------------+
                     |    | self.use\_fp8 ?  |
                     |    +--+---------+----+
                     |       | No      | Yes
                     |       |         |
                     |       |          
                     |       |   +--------------+
                     |       |   | has float8? |
                     |       |   +--+-------+---+
                     |       |      | No    | Yes
                     |       |      |       |
                     |                       
                     |    +----------------------+
                     +--- |  SDPA fallback       |
                          +----------------------+
                                    |
                                     
                          +----------------------+
                          |  Output              |
                          +----------------------+
\end{lstlisting}

\subsection{Detection Logic}

\subsubsection{FlashAttention Import}

\begin{lstlisting}
try:
    from flash\_attn\_interface import flash\_attn\_func
    FLASH\_ATTN\_AVAILABLE = True
except ImportError:
    try:
        from flash\_attn import flash\_attn\_func
        FLASH\_ATTN\_AVAILABLE = True
    except ImportError:
        FLASH\_ATTN\_AVAILABLE = False
\end{lstlisting}

The first import path is used by the Hopper-optimized \texttt{flash-attn-3} package; 
the second by the community \texttt{flash-attn} package.

\subsubsection{FP8 Path Condition}

At each forward, we check four conditions:

\begin{lstlisting}
use\_fp8\_path = (
    FLASH\_ATTN\_AVAILABLE
    and self.use\_fp8
    and q.is\_cuda
    and hasattr(torch, "float8\_e4m3fn")
)
\end{lstlisting}

\textbf{Note.} We do \textbf{not} check CUDA capability directly; if the hardware 
does not support FP8, \texttt{flash\_attn\_func} will raise a runtime error and we 
catch it, falling back to SDPA.

\subsection{FP8 Cast}

When the FP8 path is taken:

\begin{lstlisting}
q\_fa = q.transpose(1, 2).contiguous().to(torch.float8\_e4m3fn)
k\_fa = k.transpose(1, 2).contiguous().to(torch.float8\_e4m3fn)
v\_fa = v.transpose(1, 2).contiguous().to(torch.float8\_e4m3fn)
out = flash\_attn\_func(q\_fa, k\_fa, v\_fa, causal=causal\_flag).to(dtype)
\end{lstlisting}

\textbf{Cast order.} Q and K are cast to e4m3 (higher precision for products); 
V is cast to the same. The output is returned in FP16/BF16.

\textbf{No scaling.} Unlike TransformerEngine, we do not apply per-tensor scaling. 
This is correct for well-behaved distributions but may lose precision for 
models with large activations.

\subsection{SDPA Fallback}

If the FP8 path is not taken:

\begin{lstlisting}
if causal\_flag and attn\_mask is None:
    out\_sdpa = F.scaled\_dot\_product\_attention(
        q, k, v, attn\_mask=None, is\_causal=True, dropout\_p=self.dropout.p
    )
else:
    # Build causal mask, add padding mask
    out\_sdpa = F.scaled\_dot\_product\_attention(
        q, k, v, attn\_mask=attn\_mask, is\_causal=False, dropout\_p=self.dropout.p
    )
\end{lstlisting}

PyTorch's SDPA dispatches internally to:
\begin{itemize}
\item FlashAttention-2 (if installed and CUDA \textgreater{}= 11.8)
\item Memory-efficient attention
\item Native math

\end{itemize}
\subsection{Architecture-Specific Weight Transfer}

The \texttt{copy\_weights} method extracts Q, K, V, and output projections from four 
families:

\subsubsection{GPT-2 (Fused Conv1D)}

\begin{lstlisting}
if hasattr(src, "c\_attn"):
    w = src.c\_attn.weight.data
    if src.c\_attn.\_\_class\_\_.\_\_name\_\_ == "Conv1D":
        w = w.t()  # (in, out) $\to$ (out, in)
    q\_w, k\_w, v\_w = w.chunk(3, dim=0)
    extracted["q"] = (q\_w, ...)
    ...
\end{lstlisting}

\subsubsection{RoBERTa (Nested)}

\begin{lstlisting}
elif hasattr(src, "self") and hasattr(src.self, "query"):
    q\_w = src.self.query.weight.data
    k\_w = src.self.key.weight.data
    v\_w = src.self.value.weight.data
    if hasattr(src.output, "dense"):
        o\_w = src.output.dense.weight.data
\end{lstlisting}

\subsubsection{DistilBERT (Flat)}

\begin{lstlisting}
else:
    q\_w = src.q\_lin.weight.data
    k\_w = src.k\_lin.weight.data
    v\_w = src.v\_lin.weight.data
    o\_w = src.out\_lin.weight.data
\end{lstlisting}

\subsubsection{Qwen (GQA + RoPE)}

Qwen uses separate \texttt{q\_proj}, \texttt{k\_proj}, \texttt{v\_proj}, \texttt{o\_proj}. Standard extraction.

\subsection{Nested Encoder Bypass}

For RoBERTa and other BERT-family models, the output projection lives 
outside the self-attention module:

\begin{lstlisting}
BertAttention
+-- self: BertSelfAttention
|   +-- query
|   +-- key
|   +-- value
+-- output: BertSelfOutput
    +-- dense  $\leftarrow$ output projection
\end{lstlisting}

We set \texttt{is\_nested\_encoder = True}, which causes \texttt{forward} to return raw 
\texttt{context} (without applying \texttt{out\_proj}). The outer \texttt{output.dense} applies 
the projection.

For DistilBERT, \texttt{out\_lin} is inside the attention module, so we set 
\texttt{is\_nested\_encoder = False}.

\subsection{LoRA Adaptation}

Each projection is wrapped in \texttt{LoRALinear}:

\begin{lstlisting}
class LoRALinear(nn.Module):
    def \_\_init\_\_(self, in\_features, out\_features, ..., lora\_rank=16):
        self.weight = nn.Parameter(...)      # frozen after copy
        self.lora\_A = nn.Parameter(...)      # trainable
        self.lora\_B = nn.Parameter(...)      # trainable, initialized to 0
    
    def forward(self, x):
        out = F.linear(x, self.weight, self.bias)
        if self.lora\_rank \textgreater{} 0:
            out += F.linear(F.linear(x, self.lora\_A), self.lora\_B) * self.lora\_scaling
        return out
\end{lstlisting}

\subsection{Complexity Analysis}

\begin{table}[H]
\centering
\small
\begin{tabular}{lll}
\toprule
\textbf{Component} & \textbf{Time} & \textbf{Memory} \\
\midrule
Q/K/V projection & $\mathcal{O}(BLd^2)$ & $\mathcal{O}(BLd)$ \\
FlashAttention-3 FP8 & $\mathcal{O}(BL^2 d / S)$ & $\mathcal{O}(BLd + S^2)$ \\
SDPA (fallback) & $\mathcal{O}(BL^2 d)$ & $\mathcal{O}(BLd)$ \\
Output projection & $\mathcal{O}(BLd^2)$ & $\mathcal{O}(BLd)$ \\
\bottomrule
\end{tabular}
\end{table}

\textbf{FP8 path} uses ~$S^2$ SRAM per tile and reads Q/K/V once from HBM.

\textbf{SDPA path} uses memory-efficient attention with the same asymptotics but 
without the FP8 speedup.

\subsection{Parameter Analysis}

Our implementation \textbf{does not reduce parameter count}. The LoRA adapters 
add:

\[
\Delta p = r \cdot (d\_\text{in} + d\_\text{out})
\tag{4.1}
\]

per layer. For $r = 16$ and $d = 768$, this is 24,576 params per projection.

\textbf{Total overhead (DistilBERT):} 4 projections $\times$ 6 layers $\times$ 24,576 = 
\textbf{589,824 params} (+0.88\%).

\textbf{This is not a compression method.}

\subsection{Algorithm}

\textbf{Algorithm 1:} FlashAttention3FP8 forward pass
\begin{lstlisting}
Input:  X $\in$ R^{B$\times$L$\times$d}, attention\_mask
Output: Y $\in$ R^{B$\times$L$\times$d}

 1: Q $\leftarrow$ X W\_Q;  K $\leftarrow$ X W\_K;  V $\leftarrow$ X W\_V      With LoRA
 2: if RoPE: (Q, K) $\leftarrow$ Rotary(Q, K)
 3: if GQA:  (K, V) $\leftarrow$ repeat\_kv(K, V)
 4:
 5: if FLASH\_ATTN\_AVAILABLE AND use\_fp8 AND q.is\_cuda AND has\_float8:
 6:     q8, k8, v8 $\leftarrow$ cast\_to\_fp8(Q, K, V)
 7:     try:
 8:         out $\leftarrow$ flash\_attn\_func(q8, k8, v8, causal)
 9:     except:
10:         out $\leftarrow$ NULL    Fallback
11: else:
12:     out $\leftarrow$ NULL
13:
14: if out is NULL:
15:     out $\leftarrow$ SDPA(Q, K, V, mask, causal)
16:
17: if not is\_nested\_encoder:
18:     out $\leftarrow$ out\_proj(out)
19: return out
\end{lstlisting}

